% ═══════════════════════════════════════════════════════════════════════════
%  THESIS CONTENT OUTLINE
%  Augmented Reality-Based Interior Design: A Mobile Application for
%  Parametric Room Scanning and Interactive 3D Scene Reconstruction
%  Author: Mara Ielciu
% ═══════════════════════════════════════════════════════════════════════════
%
% This file describes the structure and content of the thesis. It is not
% included in the main document; it serves as a roadmap for the writing.
%
% ---------------------------------------------------------------------------
% FRONT MATTER
% ---------------------------------------------------------------------------
% - Title, author, supervisor, abstract (English)
% - Table of contents
%
% ---------------------------------------------------------------------------
% CONTEXT AND RELATED WORK (section)
% ---------------------------------------------------------------------------
% - Indoor scene understanding (ScanNet, semantic 3D)
% - AR for interior design (Kan et al., recommender, visualisation)
% - UX in mobile AR interior design (Kandil et al., SUS, HARUS)
%
% ---------------------------------------------------------------------------
% CHAPTER 1 — INTRODUCTION
% ---------------------------------------------------------------------------
% - Problem statement and goals
% - Scope of the thesis
% - Document structure (how the rest of the thesis is organised)
%
% ---------------------------------------------------------------------------
% CHAPTER 2 — THEORETICAL FOUNDATIONS I: CAPTURE AND RECONSTRUCTION
% ---------------------------------------------------------------------------
% 2.1  Motivation for the Work
%      - Why AR-based interior design; why parametric room scanning;
%        need for editable, semantic 3D; relation to prior work.
% 2.2  Photogrammetry
%      - Definition: 3D from images; overlap, parallax, triangulation.
%      - Pipeline overview: feature detection, matching, SfM, MVS,
%        surface reconstruction, texturing.
% 2.3  Object Capture and Room Plan
%      - Object Capture: Apple framework; guided capture; photogrammetry
%        pipeline on-device; USDZ output.
%      - Room Plan: parametric room scanning; walls, doors, windows,
%        furniture; CapturedRoom; relation to ARKit and LiDAR.
%
% ---------------------------------------------------------------------------
% CHAPTER 3 — THEORETICAL FOUNDATIONS II: 3D GAUSSIAN SPLATTING
% ---------------------------------------------------------------------------
% 3.1  What is 3D Gaussian Splatting
%      - Alternative to mesh-based reconstruction; primitives as 3D Gaussians.
% 3.2  Representation: 3D Gaussians
%      - Position, covariance (scale + rotation), opacity, colour (SH or RGB).
% 3.3  Splatting: Projection and Rasterisation
%      - Project 3D to 2D; sort by depth; alpha-blend.
% 3.4  Training and Relation to Object Capture
%      - Differentiable splatting; training from multi-view images;
%        comparison with mesh-based Object Capture in the application.
%
% ---------------------------------------------------------------------------
% CHAPTER 4 — IMPLEMENTATION / APPLICATION (to be written)
% ---------------------------------------------------------------------------
% - System architecture (MVVM, SwiftUI, RealityKit)
% - Room scanning flow (RoomPlan), room editor, furniture replacement
% - Object Capture integration, 3DGS viewing pipeline
% - Data persistence, concurrency
%
% ---------------------------------------------------------------------------
% CHAPTER 6 — CONCLUSIONS
% ---------------------------------------------------------------------------
% - Summary of contributions
% - Limitations and future work
%
% ---------------------------------------------------------------------------
% BIBLIOGRAPHY
% ---------------------------------------------------------------------------
% - references.bib
